In this section we prove the independence of the Hamiltonians associated with the Lax matrix.

\begin{proposition}
    In the case $\mathfrak{sl}(r+1)$, the functions $H_2,\dots,H_{r+1}$, with
    \begin{equation}
        H_k=\Tr\left(L^k\right),
    \end{equation}
    are functionally independent on a nonempty open subset of phase space.
\end{proposition}

\begin{proof}
    In the case $\mathfrak{sl}(r+1)$ one has $\Tr L = 0$, hence $H_1=0$ identically.
    Therefore we restrict to the functions $H_2,\dots,H_{r+1}$, which provide $r$ functions.

    The variables $p_i$ satisfy the constraint $\sum_{i=1}^{r+1} p_i = 0$, so we may regard $p_1,\dots,p_r$ as independent variables.

    For the Toda Lax matrix, one has
    \begin{equation}
        H_k = \sum_{i=1}^{r+1} p_i^k + R_k(p,q),
    \end{equation}
    where $R_k(p,q)$ has degree in $p$ strictly smaller than $k$.

    We first prove that the functions
    \begin{equation}
        F_k(p_1,\dots,p_r) = \sum_{i=1}^{r+1} p_i^k,
        \qquad k=2,\dots,r+1,
    \end{equation}
    are functionally independent on the hyperplane $\sum_i p_i = 0$.

    Writing
    \begin{equation}
        p_{r+1} = -\sum_{i=1}^r p_i,
    \end{equation}
    we compute the Jacobian of the functions $F_k$ with respect to the independent variables $p_1,\dots,p_r$:
    \begin{equation}
        \pdv{F_k}{p_j}
        = k p_j^{k-1} + k p_{r+1}^{\,k-1}\pdv{p_{r+1}}{p_j}
        = k\left(p_j^{k-1} - p_{r+1}^{\,k-1}\right),
        \qquad j=1,\dots,r.
    \end{equation}

    Thus the Jacobian matrix $J_F$ is
    \begin{equation}
        (J_F)_{kj} = k\left(p_j^{k-1} - p_{r+1}^{\,k-1}\right),
        \qquad k=2,\dots,r+1,\quad j=1,\dots,r.
    \end{equation}
    
    Up to the nonzero factors $k$, this matrix has the structure of a Vandermonde-type matrix built from the variables $p_1,\dots,p_{r+1}$.
    A direct computation shows that its determinant is proportional to
    \begin{equation}
        \prod_{1 \le i < j \le r+1} (p_i - p_j),
    \end{equation}
    and is therefore nonzero whenever the $p_i$ are pairwise distinct.
    Hence the Jacobian $J_F$ is non-degenerate on a nonempty open subset of the hyperplane $\sum_i p_i=0$, and the functions $F_2,\dots,F_{r+1}$ are functionally independent there.

    We now return to the functions $H_k$. Their Jacobian satisfies
    \begin{equation}
        J_H = J_F + J_R,
        \qquad
        (J_R)_{kj} = \pdv{R_k(p,q)}{p_j}.
    \end{equation}
    Since $\deg_p R_k < k$, we have
    \begin{equation}
        \deg_p \pdv{R_k}{p_j} < k-1,
    \end{equation}
    so every entry of $J_R$ has degree in $p$ strictly smaller than the corresponding entry of $J_F$.

    It follows that $\det J_H$ has the same leading homogeneous term, as a polynomial in $p$, as $\det J_F$.
    In particular, since $\det J_F$ is not identically zero, neither is $\det J_H$.

    Hence $\det J_H \neq 0$ on a nonempty open subset of phase space.
    Therefore the functions $H_2,\dots,H_{r+1}$ are functionally independent on a nonempty open subset of phase space.
\end{proof}
